
% #######################################
% ########### FILL THESE IN #############
% #######################################
% #######################################
\def\mytitle{COMBINATIONAL CIRCUITS}
\def\mykeywords{}
\def\myauthor{SRIKANTH REDDY S}
\def\contact{ssrikanthreddy03@gmail.com}
\def\mymodule{ Future Wireless Communication(FWC220107)}
% #######################################
% #### YOU DON'T NEED TO TOUCH BELOW ####
% #######################################
\documentclass[10pt, a4paper]{article}
\usepackage[a4paper,outer=1.5cm,inner=1.5cm,top=1.75cm,bottom=1.5cm]{geometry}
\twocolumn
\usepackage{circuitikz}
\usepackage{graphicx}
\graphicspath{{./images/}}
%colour our links, remove weird boxes
\usepackage[colorlinks,linkcolor={black},citecolor={blue!80!black},urlcolor={blue!80!black}]{hyperref}
%Stop indentation on new paragraphs
\usepackage[parfill]{parskip}
%% Arial-like font
\usepackage{lmodern}
\renewcommand*\familydefault{\sfdefault}
%Napier logo top right
\usepackage{watermark}
%Lorem Ipusm dolor please don't leave any in you final report ;)
\usepackage{karnaugh-map} 
\usepackage{tabularx}
\usepackage{lipsum}
\usepackage{xcolor}
\usepackage{listings}
%give us the Capital H that we all know and love
\usepackage{float}
%tone down the line spacing after section titles
\usepackage{titlesec}
%Cool maths printing
\usepackage{amsmath}
%PseudoCode
\usepackage{algorithm2e}
\usepackage{circuitikz}
\usetikzlibrary{calc}

\titlespacing{\subsection}{0pt}{\parskip}{-3pt}
\titlespacing{\subsubsection}{0pt}{\parskip}{-\parskip}
\titlespacing{\paragraph}{0pt}{\parskip}{\parskip}
\newcommand{\figuremacro}[5]{
    \begin{figure}[#1]
        \centering
        \includegraphics[width=#5\columnwidth]{#2}
        \caption[#3]{\textbf{#3}#4}
        \label{fig:#2}
    \end{figure}
}


 \lstset{
frame=single, 
breaklines=true,
columns=fullflexible
}

\thiswatermark{\centering \put(1,-110){\includegraphics[scale=0.1]{IIT_logo.png}} }
\title{\mytitle}
\author{\myauthor\hspace{1em}\\\contact\\IITH\hspace{0.5em}-\hspace{0.5em}\mymodule}
\date{}
\hypersetup{pdfauthor=\myauthor,pdftitle=\mytitle,pdfkeywords=\mykeywords}
\sloppy
% #######################################
% ########### START FROM HERE ###########
% #######################################
\begin{document}
 \maketitle
 \begin{abstract}
     %Replace the lipsum command with actual text 
  Q(19)2016 GATE: In the Digital Circuit given below,F is. 
 \end{abstract}
   
 
\begin{tikzpicture}
 

 
% Logic ports
\node[nand port] (a) at (2,1){};
\node[nand port] (b) at (2,4){};
\node[nand port] (c) at (4,0){};
\node[nand port] (d) at (6,3){};

 
% Connection

 
\draw (a.in 2) -| (b.in 2);
\draw (b.out) -| (d.in 1);
 
\draw (a.out) -|  (c.in 1);
\draw (c.out) -| (d.in 2);
\draw (d.out) -- ++(1,0) node[near end,above]{F};
 
\draw (b.in 1) -- ++(-1.5,0)node[left](In1){X};
\draw (b.in 2) -- ++(-1.5,0)node[left](In3){Y};
\draw (c.in 2) -- ++(-1.5,0)node[left](In3){Z};
% Jump crossing element
1
\end{tikzpicture}

  (1) XY+YZ'        (2)XY+Y'Z
  (3) X'Y'+YZ'       (4) XZ'+Y'
 \section{Components}
 
       \begin{tabularx}{0.4\textwidth} { 
  | >{\centering\arraybackslash}X 
  | >{\centering\arraybackslash}X 
  | >{\centering\arraybackslash}X | }
\hline
 \textbf{Components}& \textbf{Values} & \textbf{Quantity}\\
\hline
Arduino & UNO & 1 \\  
\hline
JumperWires& M-M & 6 \\ 
\hline
Breadboard &  & 1 \\
\hline

LED & & 1 \\
\hline
\end{tabularx}

    




 \section{Logic}
 The circuit takes 3-bit number from (0-7) as input X,Y,Z and produces the F as output according to the logic given in table 1.
\begin{table}[htbp]
 \begin{center}
    \begin{tabular}{|l|c|c|c|c|c|c|c|c|} \hline 
  \textbf{X}& \textbf{Y} & \textbf{Z} &\textbf{F(XY+!YZ)} \\
 \hline
 0&0&0&0\\ \hline
0&0&1&0 \\ \hline
0&1&0&1\\ \hline
0&1&1&0  \\ \hline
1&0&0&0\\ \hline
1&0&1&0\\ \hline
1&1&0&1\\ \hline
1&1&1&1\\ \hline
\end{tabular}   
\end{center}
\caption{\label{•}abel{table:dummytable} }
\end{table}
 
   
  

    
\section{karnaugh map}

Using the boolean logic output F can be expressed in terms of the inputs X,Y,Z with the help of the following Kmap.
\\
\\
\\
     \begin{karnaugh-map}[4][2][1][$YZ$][$X$]
        \minterms{2,6,7}
        \maxterms{2,6}
        \maxterms{0,1,3,4,5}
        \implicant{2}{6}
        \implicant{7}{6}
    \end{karnaugh-map}
%\end{center}
%\end(document)
\begin{center}
Fig. 2
\end{center}

 
     
    \section{Hardware Connection}


    
    \begin{table}[htbp]
 \begin{center}
    \begin{tabular}{|l|c|c|c|c|c|c|} \hline 
  \textbf{Arduino}& \textbf{2} & \textbf{3}&\textbf{4}&\textbf{5} &\textbf{GND} &\textbf{Vcc}\\
   \hline
 breadboard& F & X & Y & Z & +ve & -ve\\ \hline
led & - & - & - & - & - \\ \hline
\end{tabular}   
\end{center}
\caption{\label{table:dummytable} }
\end{table}
Give the connections as per Table 2. For taking the inputs connect 5V of arduino to +ve line of bread board to consider it as logic 'HIGH'.connect GND pin of arduino to -ve line of bread board to consider it as logic 'LOW'.
\\
\\
For example if the inputs X,Y,Z are connected 0,1,0 respectively the output should be 1 i.e., the LED connected to the GND LED should glow.
\\
\\
In the another case if we connect the inputs X,Y,Z to 0,0,0 respectively the output should be 0 i.e., the LED connected to GND pin should turn off

The circuit implementation of the above function is given in figure 1.




  \section{Procedure}
  1.Connect the arduino to the USB port of computer
  \\
  \\2.Download the follwing code
  \\
  \begin{lstlisting}
    https://github.com/srikanth9515/FWC/tree/main/IDE%20ASSIGMENT/CODES
  \end{lstlisting}
  
  3.Upload the code into the arduino board.
  
  \begin{lstlisting}
    pio run -t upload
  \end{lstlisting}
  
  4.The output '1' is represented as the state:'LED ON' and '0' is represented as the state 'LED OFF'
   



    

    
 


\end{document}

